\section{Results}
\label{sec:results}

\subsection{Most selected contexts need no finer modeled dates}

Table~\ref{tab:mechanism} reports exact context decisions under the frozen title-aware ranking.
Among 417 template contexts containing modeled evidence, 313 remain unchanged across all permitted timelines.
The corresponding human group has 79 stable contexts among 90.
Their stability rates are 75.06\% and 87.78\%, respectively.
The remaining contexts contain selected support decisions that date refinement must resolve.
\begin{table}[t]
\centering\small
\begin{tabular}{@{}lrr@{}}
\toprule
Certificate cohort & Template & Human\\
\midrule
Questions & 2,348 & 830\\
Parsed query periods & 2,296 & 813\\
Contexts with modeled evidence & 417 & 90\\
Certified stable contexts & 313 & 79\\
Unstable contexts & 104 & 11\\
Stable within modeled group & 75.06\% & 87.78\%\\
\bottomrule
\end{tabular}
\caption{Stability for the frozen title-aware ranking. The modeled group contains parsed questions with at least one modeled excerpt in the possible context. All-fallback contexts have a separate denominator.}
\label{tab:mechanism}
\end{table}


The adapter extracts 792 development claims and 744 test claims.
They occupy 2.37\% and 2.16\% of indexed source units, respectively.
Explicit end events number 445 and 384.
The confirmation cases exercise uncertain starts and explicit ends.
All three recovered succession pairs belong to calibration articles.
Appendix~\ref{app:pipeline} reports the independent source audit and retained extraction judgments.

Parsed all-fallback contexts number 1,879 and 723 across the two tracks.
Their evidence stays fixed under the shared fallback policy.
We keep them outside the modeled-context stability rates.

\subsection{The budget exposes consequential uncertainty}

The first unresolved rank gives the largest stable context prefix.
Table~\ref{tab:budget} varies only the excerpt budget on each fixed modeled cohort.
At one excerpt, stability reaches 89.9\% for template questions and 96.7\% for human questions.
At twenty excerpts, both rates approach 73\%.
A system can therefore retain a short stable context while exposing the extra evidence requiring date verification.
\begin{table}[t]
\centering\small
\begin{tabular}{@{}lrrrrr@{}}
\toprule
Excerpt budget $k$ & 1 & 3 & 5 & 10 & 20\\
\midrule
Template ($n=417$) & 89.9 & 77.2 & 75.1 & 74.1 & 73.1\\
Human ($n=90$) & 96.7 & 90.0 & 87.8 & 82.2 & 73.3\\
\bottomrule
\end{tabular}
\caption{Stable contexts (\%) as the unit budget grows. Each row retains the same modeled cohort selected at $k=5$. Selection precedes word truncation.}
\label{tab:budget}
\end{table}


The lazy selector examines only the candidates needed to fill the context.
On modeled template contexts, it examines 5.26 candidates and performs 2.92 temporal tests on average.
The human averages are 5.17 candidates and 2.70 tests.
The full masks would evaluate both support predicates for all 792 or 744 claims.
All 3,109 parsed queries match the full-mask selections exactly.
These counts isolate selection after the relevance ranking has been computed.

\subsection{Unstable contexts return replayable alternatives}

Every unstable context receives two complete assignments to the modeled event dates.
Independent replay verifies all 230 assignments and all 115 changes in rendered source excerpts.
Across both tracks, 75 exclusions retire the pivot by the query start.
Forty move its start beyond the query end.
One hundred refinement frontiers contain one claim, fourteen contain two, and one contains three.
The interface returns these current support decisions with their source spans.

\paragraph{A month-level departure changes the context.}
A human question about Helen Clark specifies the period 1989 through August 1989.
The source dates her conservation tenure from August 1987 until January 1989.
An end on January 1 removes that excerpt from the requested period.
An end on January 31 retains it.
Both dates respect the source's month-level precision.
The command-line interface reproduces both contexts from the frozen source records.
More precise evidence needs to settle this boundary; the system does not require every historical date to become exact.

\subsection{Stronger retrieval and exact certificates compose}

Table~\ref{tab:matched} compares shown-span recall under identical source units and budgets.
The pilot-selected title-aware ranking reaches 66.30\% on template questions and 57.39\% on human questions.
It improves over original BM25 with recency by 13.32 and 5.00 percentage points.
The paired intervals are $[10.80,15.80]$ and $[0.59,9.21]$ points.
\begin{table}[t]
\centering\small
\begin{tabular}{@{}lrr@{}}
\toprule
Method & Template & Human\\
\midrule
Original BM25 & 37.84 & 43.90\\
Original BM25 + latest year & 52.98 & 52.39\\
Title-aware ranking & 66.30 & 57.39\\
\quad + earliest completion & \textbf{66.84} & 57.87\\
\quad + midpoint completion & 66.56 & 57.75\\
\quad + latest completion & 66.48 & 57.75\\
\quad + interval control & 66.73 & 57.75\\
\quad + possible support & 66.73 & 57.75\\
\quad + guaranteed support & 66.62 & 57.87\\
Title-aware + latest year & 65.33 & \textbf{63.59}\\
\bottomrule
\end{tabular}
\caption{Mean question-level annotated span recall (\%) within five excerpts and 768 source-text words. Temporal controls share the title-aware ranking. Latest-year controls change only ranking.}
\label{tab:matched}
\end{table}


Adding recency to the title-aware ranking reaches 63.59\% on human questions.
Its improvement over the original recency control is 11.20 points, with interval $[7.52,15.19]$.
Possible support and the interval control select identical excerpts on these tracks.
The same evidence pipeline therefore combines improved retrieval with exact uncertainty certificates.

A secondary analysis applies the certificate to all six prespecified ranking variants.
It reads no answer labels and preserves the frozen main ranking.
Under the strongest human retrieval control, 146 of 166 modeled contexts remain stable.
Across variants, mean temporal tests range from 2.55 to 3.00 per modeled context.
The appendix reports variant-specific coverage and fixed-cohort comparisons.

\subsection{Answer generation under matched contexts}
Table~\ref{tab:reader_main} compares generated answers under the same evidence budgets.
The selected ranking improves over original BM25 with recency by 6.67 strict F1 points.
Its paired interval is $[1.59,13.56]$.
Possible and guaranteed support produce identical answer sets on the primary sample.
% Generated only from complete frozen-scoring outputs.
\begin{table}[t]
\centering\small
\begin{tabular}{@{}lrr@{}}
\toprule
Method & Primary & Diagnostic\\
\midrule
Original BM25 & 18.33 & 9.09\\
Original BM25 + year & 15.00 & 0.00\\
Title-aware ranking & 21.67 & 0.00\\
\quad + earliest completion & 20.00 & 24.24\\
\quad + midpoint completion & 20.00 & 6.06\\
\quad + latest completion & 20.00 & 6.06\\
\quad + interval control & 20.00 & 6.06\\
\quad + possible support & 20.00 & 6.06\\
\quad + guaranteed support & 20.00 & 24.24\\
Title-aware + year & 21.67 & 0.00\\
\bottomrule
\end{tabular}
\caption{Strict answer-set F1 (\%). The primary sample has 60 questions. The diagnostic selects eleven questions with differing possible and guaranteed contexts.}
\label{tab:reader_main}
\end{table}


Their answer sets differ on six diagnostic questions, including three with two valid responses.
Earliest and latest date completion yield the same change counts.
The diagnostic connects context sensitivity to observed answer changes.
Appendix~\ref{app:reader} reports exact match, paired intervals, and response-format failures.

